\section{Direct strict vanishing for the conductor-twenty character}
\label{sec:cubic-psi20-strict-vanishing}

Two explicit units complete the original-field calculation for
\(\psi_{20}=\chi_{-20}\).  The class inputs are existing certified
degree-six computations.  The unit checks below use exact rational
arithmetic and do not require a complete fundamental-unit system.
The resulting strict vanishing is independent of the host and
model transfer in AC-G124.

\paragraph{Original field, coefficients and local conditions.}
Put \(p=269\), \(K=\mathbb Q(\sqrt{-3})\), \(t=\sqrt[3]{269}\),
and let \(\chi:G_K\to\mu_3\) be either generator character of
\(K(t)/K\).  Let \(c_p\) be the actual action on roots of unity,
so \(c_p(\operatorname{Fr}^{\mathrm{geom}}_v)=q_v^{-1}\), and
let \(\psi_{20}\) be the primitive quadratic character of
\(\mathbb Q(\sqrt{-5})\), of conductor \(20\).
For \(j=1,2\) and \(e\ge1\), write
\[
                 a_{j,e}=\chi^{-j}\psi_{20}^{-1}c_p^{-e}|_{G_K}.
 \tag{G126.1}
\]
The strict group for any discrete coefficient module \(V\) means
exactly
\[
 \operatorname{Sel}_{\mathrm{str}}(K,V)
  =\ker\left(H^1_{\mathrm{cts}}(G_K,V)\longrightarrow
     \prod_{v\mid p}H^1_{\mathrm{cts}}(G_{K_v},V)
       \times\prod_{v\nmid p,\infty}H^1_{\mathrm{cts}}(I_v,V)\right).
 \tag{G126.2}
\]
In particular the coefficient-prime condition is zero in the
full local cohomology group.  All finite away inertia groups,
including those over \(2,3,5\), are retained.

\begin{theorem}[AC-G126-residual: the two conductor-twenty residual groups]
\label{AC-G126-residual}
For every finite field \(k\) of characteristic \(269\) containing
the character values, both cubic labels, and
\[
                         e\equiv267\text{ or }133\pmod {268},
 \tag{G126.3}
\]
one has
\[
              \operatorname{Sel}_{\mathrm{str}}
                      (K,k(\bar a_{j,e}))=0.
 \tag{G126.4}
\]
The finite arithmetic inputs are the certified class group
of \(\mathbb Q(t,\sqrt{15})\), the certified class-group
quotient for \(\mathbb Q(t,\sqrt{4035})\), and the two
single-unit calculations below.  There is no automorphic
endpoint, auxiliary host or deformation-model premise.
\end{theorem}

\begin{proof}
\proofstep{1. The inverse character and its exact splitting fields.}
Set \(\eta_{j,e}=\bar a_{j,e}\bar c_p^{-1}\).
For \(e\equiv267\), this is
\(\bar\chi^{-j}\psi_{20}\).  For \(e\equiv133\), it is
\(\bar\chi^{-j}\psi_{20}\bar c_p^{134}\).
The last quadratic factor cuts out \(\mathbb Q(\sqrt{269})\),
with the sign proved by the quadratic Gauss-sum calculation
in AC-G115.  Multiplication by the character of \(K/\mathbb Q\)
does not change a restricted character on \(G_K\).  Consequently
the quadratic factors of \(\eta_{j,e}\) cut out, over \(K\),
the real fields of radicands \(15\) and \(4035=15\cdot269\),
respectively.  In the second case the rational product before
restriction has field \(\mathbb Q(\sqrt{-1345})\).

For \(d\in\{15,4035\}\) put
\[
 s_d=\sqrt d,\qquad L_d=\mathbb Q(t,s_d),\qquad M_d=K(t,s_d).
 \tag{G126.5}
\]
The splitting field of \(X^3-269\) has Galois group \(S_3\)
and unique quadratic subfield \(K\).  Neither real quadratic
field in (G126.5) is \(K\).  Thus
\[
 \operatorname{Gal}(M_d/\mathbb Q)=S_3\times C_2,\qquad
 \Delta_d=\operatorname{Gal}(M_d/K)=C_3\times C_2.
 \tag{G126.6}
\]
The nontrivial element \(b\) of the second factor fixes \(K(t)\)
and sends \(s_d\) to \(-s_d\).  The character \(\eta_{j,e}\)
has nontrivial cubic part and \(b\)-value \(-1\); its kernel
over \(K\) is exactly \(G_{M_d}\).
The Kummer character to select is \(\eta_{j,e}^{-1}\), not
\(\eta_{j,e}\).

Here is the precise descent.  Let \(S_d\) be all primes of
\(M_d\) above \(p\), and put
\[
 \begin{split}
 V_{d,S_d}
  &=\{x\in M_d^\times:
       p\mid\operatorname{ord}_w(x)\text{ for every }w\nmid p\}
                                                 /M_d^{\times p},\\
 \mathcal L_d&=\bigoplus_{w\mid p}M_{d,w}^\times/M_{d,w}^{\times p}.
 \end{split}
 \tag{G126.7}
\]
Over \(M_d\) the coefficient is \(\mu_p\otimes_{\mathbb F_p}k\).
The natural action on its \(H^1\), when viewed as the
restriction of \(k(\bar a_{j,e})\), is the Kummer action
multiplied by \(\eta_{j,e}\).  Prime-to-\(p\)
inflation--restriction therefore identifies
\[
 \operatorname{Sel}_{\mathrm{str}}(K,k(\bar a_{j,e}))
   \simeq e_{\eta_{j,e}^{-1}}
       \ker\left(V_{d,S_d}\otimes_{\mathbb F_p}k
                       \longrightarrow\mathcal L_d\otimes_{\mathbb F_p}k\right).
 \tag{G126.8}
\]
For completeness, the finite global quotient and each local
decomposition or inertia quotient have order dividing six.
Their positive cohomology with characteristic-\(p\) coefficients
vanishes by averaging.  Thus restriction detects every zero
local class, including at the ramified away places.
The Kummer sequence identifies unramifiedness away from \(p\)
with the valuation condition in (G126.7); at \(p\), zero
means an actual local \(p\)th power.
This proves the local-kernel identification in (G126.8)
and does not require \(\mu_p\subset M_d\).

\proofstep{2. The two exact class certificates and the target class part.}
For \(d=15\), put \(y=s_{15}-t\).  Its polynomial is
\[
 P_{15}(Y)=Y^6-45Y^4+538Y^3+675Y^2+24210Y+68986.
 \tag{G126.9}
\]
Indeed
\[
 s_{15}=\frac{y^3+45y+269}{3y^2+15},
\]
and the denominator is nonzero at the chosen real embedding.
Squaring this identity gives (G126.9); it also proves that
\(\mathbb Q(y)=L_{15}\).
The retained PARI 2.19.1 run for this exact polynomial has
\[
 \begin{gathered}
  \texttt{PROVISIONAL\_CLASS=[14,2]},\qquad
  \texttt{PROVISIONAL\_CLASSNUMBER=28},\\
  \texttt{CERTIFY=1},\qquad \texttt{DONE}.
 \end{gathered}
 \tag{G126.10}
\]
The subsequent successful default
\texttt{bnfcertify(b)} certifies the displayed class data.
This is the existing
full \(D=15\) certificate from AC-G103, not a new class computation.

For \(d=4035\), put \(u=s_{4035}/t\).  Then
\[
 u^2=15t,\qquad u^3=15s_{4035},\qquad u^6=907875=15^3\cdot269.
 \tag{G126.11}
\]
These identities prove \(L_{4035}=\mathbb Q(u)\).
The polynomial \(Y^6-907875\) is Eisenstein at \(269\).
The retained successful run for that exact polynomial ends with
\[
 \begin{gathered}
 \texttt{PROVISIONAL\_CLASS=[14,2,2]},\\
 \texttt{CLASS\_QUOTIENT\_CERTIFY=1},\qquad \texttt{DONE}.
 \end{gathered}
 \tag{G126.12}
\]
Here the call is \texttt{bnfcertify(b,1)}.  The official manual's
conclusion is only that the actual class group is a quotient
of \(C_{14}\times C_2\times C_2\); that is all that is used.
In particular both (G126.10) and (G126.12) imply
\[
                         \operatorname{Cl}(L_d)[269]=0.
 \tag{G126.13}
\]
No full class equality or complete unit basis is inferred
from the second certificate.

Let \(r\) be the reflection in \(S_3\) fixing \(t\).
It also fixes \(s_d\), so \(M_d^{\langle r\rangle}=L_d\).
Extension and norm on classes satisfy \(Nj=2\) and
\(jN=1+r\).  Because two is invertible on \(p\)-primary
groups, (G126.13) gives
\[
        (\operatorname{Cl}(M_d)[p]\otimes_{\mathbb F_p}k)
                         ^{\langle r\rangle}=0.
 \tag{G126.14}
\]
A nonzero vector \(x\) in a target cubic eigenspace would
give a nonzero \(r\)-invariant vector \(x+rx\): reflection
inverts the cubic character, and its two nontrivial
eigenspaces are distinct.  This contradicts (G126.14).
Thus both target class parts vanish.
Apply the prime-to-\(p\) integral projector
\[
       e_*=\frac{1-b}{2}\left(1-\frac{1+c+c^2}{3}\right),
       \qquad c\text{ a generator of }C_3,
 \tag{G126.15}
\]
to the finite \(p\)-primary class group.  A nonzero component
would have nonzero \(p\)-torsion, already excluded.
Its \(S_d\)-class quotient consequently has zero target
part as well.  The Kummer unit/class sequence now reduces
the target in (G126.8) to the \(S_d\)-unit localization.

\proofstep{3. The unit multiplicity and the exact arithmetic certificates.}
The complex places of \(M_d\), with one orientation of
their restriction to \(K\), form the regular \(\Delta_d\)-set.
The logarithmic unit theorem therefore gives the regular
representation minus the trivial character on units.
There is no \(269\)-torsion unit, since
\([M_d:\mathbb Q]=12<268\).
Prime-to-\(p\) integral character projectors preserve these
multiplicities modulo \(p\).
Every \(p\)-decomposition group contains the full cubic
inertia subgroup.  Hence the added \(S_d\)-prime permutation
module has no nontrivial cubic character.
The target \(S_d\)-unit part is thus two-dimensional over
\(\mathbb F_{269}\), with \(b=-1\), and has one line of
each cubic character over \(k\).
It is simple as an \(\mathbb F_{269}[C_3]\)-module, because
\(X^2+X+1\) is irreducible modulo \(269\).

The following are the actual units used.
In the ordered basis \((1,s_{15},t,ts_{15},t^2,t^2s_{15})\),
write \(U_{15}=\frac13\sum_{i=0}^5 A_i\,v_i\), where
\begingroup\small
\[
\begin{array}{c|r}
i&A_i\\ \hline
0&-7452447141305778582247405089012236772853229436338540268869\\
1&-98767714312239608485109262799761206542363691423397314634286\\
2&-63374191728925364741595523089008606146388311014308038915873\\
3&764847356664912955555716258266672389006034176536447964081\\
4&-4041517437591935945088846457555068764710681191960788782457\\
5&-1372834474498644545069980633740298900426353455323026799531
\end{array}
\]
\endgroup
In \(L_{4035}=\mathbb Q(u)\), write
\(U_{4035}=\frac1{675}\sum_{i=0}^5 B_i u^i\), where
\begingroup\small
\[
\begin{array}{c|r}
i&B_i\\ \hline
0&331905049361710969896413766236528779053150\\
1&26340135983522172616179627067536995559525\\
2&-750696144683203818802019760302097822855\\
3&-348289924324029865187598167077589944875\\
4&-27639377212110558379124296193202143263\\
5&788360125024245269917546222251360909
\end{array}
\]
\endgroup
These rational vectors, rather than an unverified claim
about fundamental units, are the input of the unit certificate.

Here is the complete finite calculation specified by that
certificate.  Form the rational six-by-six multiplication
matrix \(M(U_d)\) in the displayed basis.  For \(d=15\),
reduce products by \(t^3=269\), \(s_{15}^2=15\);
for \(d=4035\), reduce by \(u^6=907875\).
The exact characteristic polynomial
\[
               P_{U_d}(Z)=\det(ZI-M(U_d))
 \tag{G126.16}
\]
has integral coefficients, is monic and reciprocal, and
has constant term one.  The retained rational-arithmetic
verifier prints all seven integer coefficients and
independently substitutes the polynomial in the defined
six-dimensional algebra, obtaining zero.
Thus \(U_d\) is integral, and the constant term one expresses
\(U_d^{-1}\) as an integral polynomial in \(U_d\).
This proves that both displayed elements are global units,
including at \(3\) and \(5\); no denominator is simply
discarded.  The complete input, output and rational verifier
are part of the numerical certificate.

Set, in \(M_d\),
\[
                  v_d=U_d/b(U_d),\qquad
                  w_d=v_d^3/N_{C_3}(v_d).
 \tag{G126.17}
\]
These are global units.  Since \(b\) commutes with \(C_3\),
one has \(b(w_d)=w_d^{-1}\) and \(N_{C_3}(w_d)=1\).
Consequently the class of \(w_d\) is in the target unit
module selected by (G126.15).

\proofstep{4. Nonzero localization at \(269\), with all completions retained.}
For \(d=15\), the Legendre symbol \((15/269)\) is \(-1\).
The unramified quadratic field \(K_p\) already contains
both square roots of \(15\).
There are two primes of \(M_{15}\) above \(p\), each with
absolute ramification index three and residue degree two.
At either chosen embedding, normalize \(v(t)=1\).
Reduction of the displayed vector gives, in
\(\mathbb F_{269}[s]/(s^2-15)\),
\[
 \begin{split}
 U_{15}&=A+Bt\pmod {t^2},\\
 A&=114+35s,\qquad B=241+80s,\qquad B/A=147s .
 \end{split}
 \tag{G126.18}
\]
In particular \(A\) has norm one.
The logarithmic linear coefficient of \(v_{15}\) is
\(147s-(-147s)=294s\).
The cubic norm has no \(t\)-linear term; its constant
term is the cube of that of \(v_{15}\).
Therefore
\[
                         w_{15}=1+75s\,t\pmod {t^2}.
 \tag{G126.19}
\]
The coefficient is nonzero.
The localization used in (G126.8) is the sum over
\emph{both} \(p\)-completions.  The involution \(b\)
exchanges them; no individual completion is incorrectly
declared \(\Delta_{15}\)-stable.

For \(d=4035\), (G126.11) is Eisenstein over \(K_p\)
as well.  Hence \(M_{4035}\) has a unique \(p\)-prime,
of absolute ramification index six and residue degree
two.  Normalize \(v(u)=1\), so \(v(p)=6\).
Now \(b(u)=-u\), and the exact rational vector gives
\[
             U_{4035}=1+93u\pmod {u^2},\qquad
             w_{4035}=1+20u\pmod {u^2}.
 \tag{G126.20}
\]
The second coefficient is \(6\cdot93\equiv20\pmod {269}\):
division by \(b(U)\) doubles the linear coefficient,
the cube multiplies it by three, and the cubic norm
has zero linear term.  This coefficient is also nonzero.

In either completion the absolute ramification index
\(a\) is \(3\) or \(6\), strictly smaller than \(p-1\).
The convergent logarithm and exponential identify
\(1+\mathfrak m\) with \(\mathfrak m\).
Principal-unit \(p\)th powers therefore have logarithm
in \(p\mathfrak m=\mathfrak m^{a+1}\).
The residue-unit group has order \(p^2-1\), prime to
\(p\); a \(p\)th root of a unit congruent to one must
itself be a principal unit.
The nonzero linear terms in (G126.19)--(20) prove that
the respective \(w_d\) is not a local \(p\)th power.

The localization of the simple two-dimensional
\(\mathbb F_{269}[C_3]\)-module is nonzero, hence
injective.  After extension to \(k\), it is injective
on both target character lines.
Together with the class calculation and (G126.8), this
proves (G126.4) for both \(e\)-residues and both cubic
labels.
\end{proof}

\begin{theorem}[AC-G126-lifts: the original coefficient character tower]
\label{AC-G126-lifts}
Let \(E/\mathbb Q_{269}\) be any finite coefficient field
containing the cubic values, with its original valuation
ring \(O\), uniformizer \(\varpi\), residue field \(k\),
and maximal ideal \(\mathfrak m_O\).
Let \(\omega_p\) be the Teichmuller lift of \(\bar c_p\).
For any continuous
\[
 \kappa:G_K\longrightarrow1+\mathfrak m_O,\qquad
 \kappa|_{I_v}=1\quad(v\nmid p),
 \tag{G126.21}
\]
integers \(e\ge1,r\), and \(e+r\equiv133\) or \(267\pmod {268}\),
put
\[
 a=\chi^{-j}\psi_{20}^{-1}\omega_p^{-r}\kappa^{-1}c_p^{-e}|_{G_K},
 \qquad \mathcal D=(E/O)(a).
 \tag{G126.22}
\]
Then
\[
                    \operatorname{Sel}_{\mathrm{str}}(K,\mathcal D)=0.
 \tag{G126.23}
\]
Its original \(O\)-length and normalized length are zero,
without \(e_O=1\), an extension of \(O\), a finite-order
condition on \(\kappa\), or a rational extension of \(\kappa\).

In particular, for every fixed \(e\ge1\) and \(m\ge1\),
exactly \(2\cdot269^{m-1}\) finite characters \(\tau\)
of conductor dividing \(269^m\) have tame exponent \(r\)
with \(e+r\equiv133\) or \(267\pmod{268}\).
Every one of these characters gives this conclusion for
\(\psi=\chi_{-20}\tau\), over every original coefficient
field containing its values.  These labels remain
distinct after restriction to \(K\).  This is not a
primitive-conductor count.
\end{theorem}

\begin{proof}
We repeat the coefficient argument with its actual new
away character; the theorem AC-G125-lifts itself fixes
\(\chi_{-4}\) and is not invoked with a changed base character.
Reduction of (G126.22) is
\[
                 \bar a=\bar\chi^{-j}\psi_{20}^{-1}
                                             \bar c_p^{-(e+r)}.
 \tag{G126.24}
\]
It is the actual residual module in (G126.4), including
all its local restriction maps.  Thus its strict
residual group is zero.

At \(p\), \(\psi_{20}\) is unramified.  Raising an
inertia element with nontrivial cubic value to the
\(268\)th power kills its rational residual cyclotomic
value and preserves its cubic value.
Since \(\bar\kappa=1\), some \(a(g)-1\) is a unit of
the original \(O\).  Hence
\(\mathcal D^{G_{K_p}}=\mathcal D^{G_K}=0\).
Away from \(p\), \(\omega_p,c_p,\kappa\) are trivial
on inertia, so the remaining character is
\(\chi^{-j}\psi_{20}^{-1}\), of order dividing six.
Its invariants on \(\mathcal D\) are either zero or
all of \(E/O\), and are \(\varpi\)-divisible.
This explicitly includes the added ramification at \(5\).

Multiplication by the actual uniformizer gives
\[
       0\longrightarrow k(\bar a)\longrightarrow
        \mathcal D\xrightarrow{\varpi}\mathcal D\longrightarrow0.
 \tag{G126.25}
\]
If \(\varpi[z]=0\), write \(\varpi z=dm\), choose
\(\varpi n=m\), and replace \(z\) by \(z-dn\).
This is a continuous cocycle with values in
\(\mathcal D[\varpi]\simeq k(\bar a)\).
If its local restriction is \(du\) in \(\mathcal D\),
then \(\varpi u\) is locally invariant.
Local invariant divisibility supplies invariant \(u_0\)
with \(\varpi u_0=\varpi u\).
Thus \(u-u_0\) lies in \(\mathcal D[\varpi]\) and has
the same coboundary.  The residual lift of a strict
class is itself strict, and is zero by (G126.4).
Globally the lift is unique because the invariant
module is zero.  We conclude that
\(\operatorname{Sel}_{\mathrm{str}}(K,\mathcal D)[\varpi]=0\).

Every continuous cocycle from compact \(G_K\) to
discrete \(E/O\) has finite image.  One power of
\(\varpi\) kills that image, so every class is
\(\varpi\)-primary torsion.
A nonzero strict class would have a nonzero
\(\varpi\)-torsion multiple, a contradiction.
This proves (G126.23) over the original \(O\).

For the finite-character assertion use the elementary
decomposition already proved in AC-G125-Dirichlet:
\((\mathbb Z/p^m\mathbb Z)^\times\) has tame character
exponent \(r\pmod {268}\) and \(p^{m-1}\) wild
characters.  Precisely the two exponents
\(r\equiv133-e,267-e\) are allowed.
Their wild parts reduce to one and are unramified
away from \(p\), so the preceding argument applies.
The intersection \(K\cap\mathbb Q(\mu_{p^m})\) is
trivial by ramification at \(3\), proving distinctness
after restriction.
Only this finite-group/counting argument is reused;
the new residual and coefficient vanishing has been
proved explicitly above for \(\psi_{20}\).
\end{proof}

\paragraph{Certificate and geometric scope.}
The \(D=15\) default-flag class certificate is the frozen
AC-G103 computation.  The \(D=4035\) input is the actual
successful flag-one job for \(Y^6-907875\), not an unused
planned-job stub.  Both raw inputs, outputs and binary
data, the official PARI 2.19.1 \texttt{bnfcertify} manual
passage, and the complete unit verifier and outputs are
preserved in the evidence manifest.
The unit proof uses only monic integral relations with
constant term one and explicit local coefficients.
No new class computation, completeness of fundamental
units, or stronger flag-one conclusion is claimed.
The Kummer sequence, logarithmic unit representation and
prime-to-\(p\) descent retain the checked source scopes
of AC-G103--AC-G115; the inverse-character and local
kernel arguments were rederived above.

Geometrically (G126.8) is the specified character part
of Kummer torsors on \(\operatorname{Spec}\mathcal O_{M_d}[1/p]\),
trivial at every coefficient-prime completion.
The finite-cover dictionary is the previously audited
SGA~1, Expose~V, sections~4 and~7; the Kummer input is
the previously checked Stacks Project tag 03PK.
Finite torsion layers of (G126.22) are finite locally
constant sheaves away their stated ramification.
An infinite-image \(\kappa\) does not make their discrete
union a finite cover.
These theorems give no automorphic endpoint, model M4,
mixed characteristic-zero Bloch--Kato vanishing,
relaxed local condition, other coefficient prime,
all-family budget, or conclusion of H04.4.
